% =====================================================================
% Reemplaza la subsección «La Fuerza Gravitacional entre Dos Cuerpos»
% (dentro de \section{Ley de Gravitación Universal de Newton}).
% Requiere: amsmath, siunitx, tikz (ya cargados en main.tex).
% Conserva las etiquetas eq:ley_gravitacion y eq:GravitationalForce.
% =====================================================================

\subsection{La ley de gravitación universal en términos del campo}
\label{subsec:gravitacion_campo}

Cerca de la superficie de la Tierra todos los cuerpos caen con la misma
aceleración, $g_0 \approx \SI{9.8}{\meter\per\second\squared}$, sin importar
su masa, y la prueba de la Luna (subsección~\ref{subsec:prueba_luna}) muestra
que esa aceleración se debilita con la distancia como $1/r^{2}$. Esto permite describir la gravitación sin hablar primero de fuerzas:
la Tierra, o cualquier cuerpo de masa $M$, modifica el espacio a su alrededor
y crea un \textbf{campo gravitacional} $\boldsymbol{g}$, cuyo valor en cada
punto es la aceleración de caída libre en ese punto.

\paragraph{Por qué el campo decae como $1/r^{2}$.}
Representamos el campo con líneas que llegan radialmente a $M$, en un número
$N$ proporcional a la masa, $N \propto M$. Donde las líneas están más
apretadas el campo es más intenso: la magnitud del campo es la
\textbf{densidad de líneas}, es decir, el número de líneas que cruzan una
superficie perpendicular a ellas dividido por su área. En el espacio vacío
las líneas no empiezan ni terminan, así que las mismas $N$ líneas cruzan
cualquier esfera centrada en $M$. Como el área de una esfera de radio $r$ es
$4\pi r^{2}$,
\begin{equation}
  g(r) = \frac{N}{4\pi r^{2}} \propto \frac{M}{r^{2}} \,.
  \label{eq:densidad_lineas}
\end{equation}
A una distancia doble el área es cuatro veces mayor y el campo se reduce a la
cuarta parte; a una distancia triple, a la novena parte
(figura~\ref{fig:lineas_campo_esferas}). Escogiendo la constante de
proporcionalidad como $N = 4\pi G M$, la magnitud del campo es
\begin{equation}
  g(r) = \frac{G M}{r^{2}} \,,
  \label{eq:campo_g}
\end{equation}
donde $G$ es la \textbf{constante de gravitación universal},
$G \approx \SI{6.674e-11}{\newton\meter\squared\per\kilo\gram\squared}$, la
misma para todos los cuerpos.

\begin{figure}[htbp]
  \centering
  \begin{tikzpicture}[
      scale=1.1,
      linea/.style={-{Stealth[length=2.5mm]}, thick, red!75!black},
      esfera/.style={thick, dashed, blue!60!black},
    ]
    % Esferas de radio r y 2r
    \draw[esfera] (0,0) circle (1.3);
    \draw[esfera] (0,0) circle (2.6);
    % Líneas de campo: llegan radialmente a M
    \foreach \ang in {0,30,...,330} {
      \draw[linea] (\ang:3.4) -- (\ang:0.45);
    }
    % Masa central
    \fill[ball color=teal!70!black] (0,0) circle (0.38);
    \node[white, font=\bfseries] at (0,0) {$M$};
    % Rótulos de las esferas
    \node[blue!60!black, fill=white, inner sep=1.5pt, font=\small]
      at (-50:1.3) {$4\pi r^{2}$};
    \node[blue!60!black, fill=white, inner sep=1.5pt, font=\small]
      at (-50:2.6) {$4\pi (2r)^{2}$};
  \end{tikzpicture}
  \caption{Las mismas $N$ líneas de campo cruzan la esfera de radio $r$ y la de
    radio $2r$, pero la segunda tiene cuatro veces más área: la densidad de
    líneas, y con ella la magnitud del campo, es cuatro veces menor.}
  \label{fig:lineas_campo_esferas}
\end{figure}

\paragraph{Forma vectorial.}
Sea $\boldsymbol{r}$ el vector de posición del punto medido desde el centro de
$M$, con magnitud $r = \lvert\boldsymbol{r}\rvert$, y sea
$\hat{\boldsymbol{r}} = \boldsymbol{r}/r$ el vector unitario que apunta desde
$M$ hacia el punto. Como las líneas de campo apuntan \emph{hacia} $M$, el
campo tiene la dirección de $-\hat{\boldsymbol{r}}$:
\begin{equation}
  \boxed{\boldsymbol{g}(\boldsymbol{r})
    = -\frac{G M}{r^{2}}\,\hat{\boldsymbol{r}}
    = -\frac{G M}{r^{3}}\,\boldsymbol{r}} \,.
  \label{eq:campo_g_vectorial}
\end{equation}

\paragraph{Calibración con la superficie de la Tierra.}
En la superficie, $r = R_T$, el campo debe ser la aceleración de caída libre:
\begin{equation}
  g_0 = \frac{G M_T}{R_T^{2}} \approx \SI{9.8}{\meter\per\second\squared} \,,
  \label{eq:calibracion_g0}
\end{equation}
de modo que, con $R_T \approx \SI{6.371e6}{\meter}$,
\begin{equation}
  G M_T = g_0 R_T^{2}
  \approx \SI{3.98e14}{\meter\cubed\per\second\squared} \,.
  \label{eq:calibracion_GM}
\end{equation}
La caída libre solo determina el producto $G M_T$; para separar $G$ de $M_T$
hace falta medir la atracción entre masas de laboratorio, como lo hizo
Cavendish en 1798.

\paragraph{Del campo a la fuerza.}
Una masa $m$ colocada en un punto donde el campo vale $\boldsymbol{g}$ cae con
aceleración $\boldsymbol{g}$. Por la ley de la fuerza, la fuerza que siente es
$\boldsymbol{F} = m\,\boldsymbol{g}$. Si $\boldsymbol{r}$ va de $M$ hasta $m$,
la fuerza que $M$ ejerce sobre $m$ es
\begin{equation}
  \boldsymbol{F}_{M \to m} = m\,\boldsymbol{g}(\boldsymbol{r})
  = -\frac{G M m}{r^{3}}\,\boldsymbol{r} \,,
  \label{eq:GravitationalForce}
\end{equation}
cuya magnitud es la \textbf{ley de gravitación universal} de Newton:
\begin{equation}
  F = \frac{G M m}{r^{2}} \,.
  \label{eq:ley_gravitacion}
\end{equation}

\paragraph{Acción y reacción.}
La masa $m$ también crea su propio campo. En la posición de $M$, el vector
que va desde $m$ hasta $M$ es $-\boldsymbol{r}$, de modo que el campo de $m$
allí es $\boldsymbol{g}_m = -G m\,(-\boldsymbol{r})/r^{3}$, y la fuerza
sobre $M$ resulta
\begin{equation}
  \boldsymbol{F}_{m \to M} = M\,\boldsymbol{g}_m
  = \frac{G M m}{r^{3}}\,\boldsymbol{r}
  = -\boldsymbol{F}_{M \to m} \,.
  \label{eq:gravitacion_accion_reaccion}
\end{equation}
Las dos fuerzas tienen la misma magnitud y sentidos opuestos, como exige la
ley de acción y reacción: la manzana atrae a la Tierra tanto como la Tierra a
la manzana, aunque la aceleración de la Tierra sea imperceptible.

\paragraph{Características de la interacción gravitacional.}
\begin{itemize}
  \item Es \textbf{directamente proporcional} a cada una de las masas: si se
    duplica una de ellas, la fuerza se duplica.
  \item Es \textbf{inversamente proporcional} al cuadrado de la distancia: si
    la distancia se duplica, la fuerza se reduce a la cuarta parte. Es una
    \textbf{ley del inverso del cuadrado}, consecuencia de que las líneas de
    campo se reparten sobre esferas de área $4\pi r^{2}$.
  \item Es siempre \textbf{atractiva} y \textbf{central}: actúa a lo largo de
    la línea que une los centros de los dos cuerpos.
  \item Es \textbf{universal}: la misma expresión describe la caída de una
    manzana, la órbita de la Luna y el movimiento de los planetas.
\end{itemize}

\paragraph{Superposición.}
Si hay varias masas $M_i$ en las posiciones $\boldsymbol{r}_i$, el campo en
el punto $\boldsymbol{r}$ es la suma vectorial de los campos de cada una:
\begin{equation}
  \boldsymbol{g}(\boldsymbol{r})
  = -\sum_{i} \frac{G M_i}{\lvert \boldsymbol{r} - \boldsymbol{r}_i \rvert^{3}}
    \,(\boldsymbol{r} - \boldsymbol{r}_i) .
  \label{eq:campo_superposicion}
\end{equation}
